\documentclass[10pt,a4paper]{amsart}
\usepackage[margin=19mm]{geometry}
\usepackage{amsmath,amssymb,mathtools}
\usepackage[scr=rsfs,cal=euler]{mathalfa}
\usepackage{xfrac}
\usepackage[unicode,hidelinks,bookmarks=false]{hyperref}
\newcommand{\ie}{i.e.\ }
\newcommand{\cf}{cf.\ }
\newcommand{\resp}{respectively\ }
\newcommand{\Subsection}{\subsection}
\providecommand{\nobreakdash}{\nobreak\hbox{-}}
\setlength{\emergencystretch}{2em}
\multlinegap0pt
\usepackage{amssymb}
\usepackage{tikz}
\newtheorem{lemma}{Lemma}
\title{Estimating the number of real zeros of linear combinations of radicals of polynomials}
\author{Gal Binyamini, Avner Kiro, Alexander Logunov, Dmitry Novikov, Dmitrii Zakharov}
\date{Source: arXiv:2609.02871v1 (CC BY 4.0)}
\begin{document}
\maketitle

\noindent\emph{Source and licence.} Estimating the number of real zeros of linear combinations of radicals of polynomials. Version arXiv:2609.02871v1 (CC BY 4.0), distributed by arXiv under the Creative Commons Attribution 4.0 International licence, http://creativecommons.org/licenses/by/4.0/, as recorded in the arXiv licence metadata for this exact version and shown on the abstract page https://arxiv.org/abs/2609.02871v1. Full licence notice: \texttt{licenses/arxiv-2609.02871-CC-BY-4.0.txt}.\par\smallskip

\noindent\textit{Editorial scope.} Counted unit: the article's lemma Lemma:main (source0.tex lines 530--536). The article's complete original proof of it is delivered with it (source0.tex lines 538--561, one \texttt{proof} environment of 24 lines). No mathematics is written, repaired or re-derived: the statement and the proof are the article's own lines, in the article's own notation, and no retained line is substituted or re-flowed.  No further source-local context is needed: every cross-reference of every retained line resolves inside the two delivered spans.  Nothing was omitted from any retained span, so the package carries no omission record.  No retained line contains a citation, so the wrapper carries no bibliography and no bibliography entry.  No retained line contains a figure environment, an image-inclusion command or drawing code, so no image is delivered and the package declares no asset dependency.  Editorial formatting only, in addition: an AMS \texttt{amsart} preamble that restates the article's own declarations and macros used by the retained lines, namely the environments \texttt{lemma} on separate counters, together with the standard packages those lines need.  The retained environments print in this document as Lemma 1 in order of appearance, whereas the article prints the same items with the numbering of its own full document.  The complete native source of the article as distributed in the arXiv e-print for this exact version is bundled unchanged as \texttt{sources/209247/source0.tex}.
\begin{lemma} \label{Lemma:main}
Let $\beta>-1/2$ and $\lambda=-\beta(\beta+1)< \frac{1}{4}$. Let $u$ be a solution to
\[
\Delta_{\mathbb{D}} u + \lambda u=\sum_{k=1}^{n}c_{k}\delta_{p_{k}} \quad \text{ in Poincare's disc } \quad \mathbb{D},
\]
where $c_{k} \in \mathbb{R}\setminus{0}$ and $p_k$ are distinct points on $D$. Assume that $u(z) = o(1-|z|^2)^{-\beta}$ as $z \to \partial \mathbb{D}$. Then $u$ has at most $n$ nodal domains in $D$.
\end{lemma}

\begin{proof}
Consider the radial positive solution $v$: $\Delta v+\lambda v=0$ on $\mathbb{D}$ such that $v(z)= v(|z|)$ and $v(0)=1$. Then the ratio $h=\frac{u}{v}$ tends to zero near the boundary of $\mathbb{D}$. 


 The zero set and nodal domains of $h$ and $u$ are the same.
We claim that each nodal domain of $h$ contains at least one of the points $p_k$.

 First, we note that the function $u$ (and h) tend to $+\infty$ near $p_k$ if $c_k<0$ and to $-\infty$ if $c_k>0$. This can be seen by considering the Mobius transformation $T$ such that $\Delta g(T(z))+\lambda g(T(z))= -\delta_{p_k}$. Then by elliptic regularity $u(z)+c_kg(T(z))$ is smooth near $p_k$ and therefore $u$ tends to $\pm\infty$ near $p_k$ depending on the sign of $c_k$.

If a nodal domain $\Omega$ of $h$ did not contain any of $p_k$, then 
 $|h|$ must have a local maximum in $\Omega$ because $h(z) \to 0$ as $z\to \partial \Omega \cup \partial \mathbb{D}$ in Euclidean metric. We claim that $|h|$ cannot have local maxina.

Assume the contrary and identify Poincare's disc with Lobachevsky upper half-plane $\mathbb{H}$.   In the $(x,y)$ coordinates on $\mathbb{H}=\mathbb{R}_{+}^{2}$ 
the functions $u$ and $v$ satisfy $$y^{2}\Delta v+\lambda v=0, \quad  y^{2}\Delta u+\lambda u=0$$
in $\Omega$, which yields that $h=u/v$ satisfies
\[
\text{div}(v^{2}\nabla h)=0 \quad \text{in } \Omega.
\]

By the maximum principle for elliptic equations in divergence form,
the function $h$ cannot have local maxima and minima unless it is identically constants. Contradiction is obtained and therefore each nodal domain of $u$ must contain at least one of the points $p_k$ and therefore the number of nodal domains is at most $n$.


\end{proof}

\end{document}
